\documentclass[11pt]{article}

\usepackage{amsmath,amssymb,amsthm,mathtools}
\usepackage[T1]{fontenc}
\usepackage{lmodern}
\usepackage[margin=1in]{geometry}
\usepackage{enumitem}
\usepackage{hyperref}
\usepackage{microtype}
\usepackage{xcolor}

\hypersetup{colorlinks=true,linkcolor=blue,citecolor=blue,urlcolor=blue}

\newtheorem{theorem}{Theorem}[section]
\newtheorem{proposition}[theorem]{Proposition}
\newtheorem{lemma}[theorem]{Lemma}
\newtheorem{corollary}[theorem]{Corollary}
\newtheorem{remark}[theorem]{Remark}
\newtheorem{definition}[theorem]{Definition}
\newtheorem{placeholder}[theorem]{Placeholder}

\title{D5 133: the front-end one-corner reduction theorem and the true remaining M1 burden}
\author{}
\date{March 2026}

\begin{document}
\maketitle

\begin{abstract}
This note rewrites the front-end odd-$m$, $d=5$ reduction in the cleaner
``one-cornered odometer'' language proposed in the rewrite plan.
It proves that the old broad front-end burden is smaller than it looked.
Once one knows that every unresolved front-end instance reduces to a single
oriented bridge-type seed and that its live continuation is the alternate-$2$
entry branch analyzed in the compact structural note \textup{098}, all later
front-end work is automatic: the normalized seed is forced to be
\([2,2,1]/[1,4,4]\), the exact trigger family is
\(H_{L1}=\{(m-1,m-1,u,2):u\neq 2\}\), and the actual branch agrees with the
candidate one-corner odometer orbit up to first exit.
In particular, the post-entry part of the old non-trigger \textup{033}
classification is already internalized by \textup{098}; it is not an
independent manuscript burden anymore.
The genuinely remaining front-end work is therefore reduced to two finite
inputs: an endpoint-word shape theorem and an initial entry-channel theorem.
\end{abstract}

\section{Scope}

The rewrite plan proposed replacing the old front-end bookkeeping
\textup{032}/\textup{033} with one cleaner theorem:

\begin{quote}
Every unresolved odd-$m$, $d=5$ front-end instance reduces, after the accepted
static repairs and symmetries, to one oriented bridge-type one-corner seed.
\end{quote}

The present note shows how far that program can already be pushed using the
current theorem-level bundle.
The key point is that the old front-end burden naturally splits into two very
small pieces:
\begin{enumerate}[label=\textup{(\arabic*)},leftmargin=2.5em]
\item a finite endpoint-word \emph{shape} theorem saying that the unresolved
      front end has one oriented bridge-type seed shape;
\item a finite \emph{entry} theorem saying that the live continuation of that
      normalized seed is the alternate-$2$ entry branch analyzed in the compact
      structural note \textup{098}.
\end{enumerate}

Everything after that point is already theorem-level.
So the note does not claim a full first-principles front-end proof from the raw
seed-search artifacts.
Instead it isolates exactly what still remains to be promoted.

\section{Imported ingredients already at theorem level}

\subsection{The oriented bridge normal form}

We use the cyclic direction alphabet
\[
\mathcal D=\{1,2,3,4\}
\]
with cyclic order
\[
1\to 2\to 3\to 4\to 1.
\]

\begin{definition}[oriented bridge-type endpoint seed]
An \emph{oriented bridge-type endpoint seed} is a pair of length-$3$ endpoint
words
\[
L=[\ell,\ell,b],\qquad R=[b,r,r]
\]
such that:
\begin{enumerate}[label=\textup{(\roman*)},leftmargin=2.5em]
\item $b$ is the shared bridge symbol;
\item $\ell$ is the repeated left neighbor of the bridge;
\item $r$ is the repeated right neighbor of the bridge;
\item after the standard orientation convention is fixed, $\ell$ is the
      clockwise neighbor of $b$ and $r$ is the counterclockwise neighbor of
      $b$.
\end{enumerate}
\end{definition}

The normal-form lemma from Note \textup{131} is the following.

\begin{proposition}[normal form of the one-corner seed]
\label{prop:normal-form}
Every oriented bridge-type endpoint seed is equivalent, under the accepted
front-end relabeling and orientation normalizations, to the unique numerical
representative
\[
L=[2,2,1],\qquad R=[1,4,4].
\]
\end{proposition}

\begin{proof}
Relabel cyclically so that the shared bridge symbol is $b=1$.
Then the oriented left repeated neighbor is the clockwise neighbor of $1$,
namely $\ell=2$, and the oriented right repeated neighbor is the
counterclockwise neighbor of $1$, namely $r=4$.
So the normalized representative is exactly
\[
L=[2,2,1],\qquad R=[1,4,4].
\]
Uniqueness is immediate because the normalization fixes the bridge symbol and
then fixes which adjacent symbol is written on each side.
\end{proof}

\subsection{The best-seed channel package}

Once the normalized seed has been fixed, the compact structural note
\textup{098} already supplies the exact best-seed channel package.
We only record the two statements that matter for front-end reduction.

\begin{proposition}[exact trigger family]
\label{prop:trigger}
On the best-seed channel, the target hole family is exactly
\[
H_{L1}=\{(q,w,u,\Theta)=(m-1,m-1,u,2):u\neq 2\}.
\]
Equivalently, the target layer is $\Theta=2$, the target current coordinates
satisfy $q=w=m-1$, and the only excluded fiber is $u=2$.
\end{proposition}

\begin{theorem}[actual first exits and pre-exit $B$-region invariance]
\label{thm:098-core}
Fix the normalized best-seed channel with source seed
\[
L=[2,2,1],\qquad R=[1,4,4],
\]
and start at the alternate-$2$ entry state
\[
E=(m-1,1,u_{\mathrm{source}},2).
\]
Then:
\begin{enumerate}[label=\textup{(\arabic*)},leftmargin=2.5em]
\item the actual active branch agrees with the candidate one-corner orbit up
      to first exit;
\item every pre-exit actual current state is $B$-labeled;
\item the first exit to $H_{L1}$ occurs at the universal family-dependent
      target
      \[
      T_{\mathrm{reg}}=(m-1,m-2,1)
      \quad\text{or}\quad
      T_{\mathrm{exc}}=(m-2,m-1,1),
      \]
      with regular families exiting at $T_{\mathrm{reg}}$ by direction $2$ and
      the exceptional family exiting at $T_{\mathrm{exc}}$ by direction $1$;
\item no patched current class is encountered before first exit.
\end{enumerate}
\end{theorem}

\begin{remark}
The last item is already part of the proof of the compact structural theorem:
it uses only the coarse coordinate signatures of the patched classes and the
monotonicity of the corridor coordinate $w$ on the candidate branch.
That point is crucial for the present note, because it means the post-entry part
of the old non-trigger classification is already internalized.
\end{remark}

\section{What is already automatic after entry}

The previous section implies that one large part of the old front-end burden is
no longer genuinely open.
Once the normalized one-corner branch is entered, the remaining local current
classes are automatically controlled by the compact \textup{098} theorem.

\begin{proposition}[post-entry non-trigger control is internalized]
\label{prop:post-entry-internalized}
Assume the normalized one-corner seed
\[
L,R]=([2,2,1],[1,4,4])
\]
has already been reduced to the alternate-$2$ entry state
\[
E=(m-1,1,u_{\mathrm{source}},2)
\]
of the best-seed channel.
Then no separate post-entry theorem is needed to rule out the old patched
current classes on the unresolved branch: before first exit to $H_{L1}$, the
actual branch never meets any patched current class at all.
In particular, the old post-entry non-trigger control of classes
\(L1,R1,L2,R2,L3,R3\) is already contained in
Theorem~\ref{thm:098-core}.
\end{proposition}

\begin{proof}
This is exactly item~(4) of Theorem~\ref{thm:098-core}.
The compact structural theorem proves that the candidate orbit avoids all patched
current classes before first exit, and then upgrades that avoidance to the
actual branch because the actual branch agrees with the candidate orbit up to
first exit.
So once one has entered the normalized one-corner branch, there is no further
independent post-entry channel-classification problem.
\end{proof}

\begin{corollary}[what is left of the old \textup{033} burden]
\label{cor:033-shrunk}
For the purposes of the main odd-$m$, $d=5$ proof spine, the old
non-trigger part of \textup{033} does \emph{not} need to be promoted in its
historical full breadth.
What remains externally necessary is only the finite \emph{initial entry}
classification saying that the live unresolved continuation of the one-corner
seed really is the alternate-$2$ entry branch handled by
Theorem~\ref{thm:098-core}.
The subsequent patched-class avoidance is already automatic.
\end{corollary}

\section{The true remaining front-end inputs}

The previous sections isolate the genuine external front-end inputs much more
sharply than the old single $M1$ placeholder.

\begin{placeholder}[endpoint-word shape theorem]
\label{ph:shape}
After the accepted tiny static repairs and front-end symmetries, every
unresolved odd-$m$, $d=5$ front-end instance reduces to a single oriented
bridge-type endpoint seed class.
\end{placeholder}

\begin{placeholder}[initial entry theorem]
\label{ph:entry}
For the normalized one-corner seed
\[
L=[2,2,1],\qquad R=[1,4,4],
\]
the live unresolved continuation is the alternate-$2$ entry branch considered
in Theorem~\ref{thm:098-core}, with entry state
\[
E=(m-1,1,u_{\mathrm{source}},2).
\]
Equivalently, every other immediate continuation out of the normalized front-end
seed is already resolved by the accepted static/direct repair package, and the
only live unresolved continuation is this alternate-$2$ branch.
\end{placeholder}

\begin{remark}
Placeholder~\ref{ph:entry} is intentionally smaller than the historical broad
wording ``promote the non-trigger part of 033.''
By Proposition~\ref{prop:post-entry-internalized}, once the alternate-$2$
entry branch is identified, all later post-entry patched-class control is
already theorem-level.
So the remaining front-end burden is not a wide defect-template classification;
it is only a finite initial-branch identification.
\end{remark}

\section{Front-end one-corner reduction theorem}

\begin{theorem}[front-end one-corner reduction]
\label{thm:one-corner-reduction}
Assume Placeholders~\ref{ph:shape} and~\ref{ph:entry}.
Then every unresolved odd-$m$, $d=5$ front-end instance reduces to the
normalized one-corner best-seed channel with seed
\[
L=[2,2,1],\qquad R=[1,4,4],
\]
entry state
\[
E=(m-1,1,u_{\mathrm{source}},2),
\]
exact trigger family
\[
H_{L1}=\{(q,w,u,\Theta)=(m-1,m-1,u,2):u\neq 2\},
\]
and universal first exits
\[
T_{\mathrm{reg}}=(m-1,m-2,1),
\qquad
T_{\mathrm{exc}}=(m-2,m-1,1).
\]
Moreover, on that branch the actual active dynamics agree with the candidate
one-corner odometer orbit up to first exit and avoid all patched current classes
before exit.
\end{theorem}

\begin{proof}
By Placeholder~\ref{ph:shape}, every unresolved front-end instance reduces,
after the accepted static repairs and symmetries, to one oriented bridge-type
endpoint seed class.
By Proposition~\ref{prop:normal-form}, that seed class has the unique normalized
representative
\[
L=[2,2,1],\qquad R=[1,4,4].
\]
By Placeholder~\ref{ph:entry}, the live unresolved continuation of this
normalized seed is the alternate-$2$ entry branch with entry state
\[
E=(m-1,1,u_{\mathrm{source}},2).
\]
Now Proposition~\ref{prop:trigger} supplies the exact trigger family
\[
H_{L1}=\{(m-1,m-1,u,2):u\neq 2\},
\]
and Theorem~\ref{thm:098-core} gives the candidate/actual agreement up to first
exit, the universal family-dependent first exits
\[
T_{\mathrm{reg}}=(m-1,m-2,1),
\qquad
T_{\mathrm{exc}}=(m-2,m-1,1),
\]
and the avoidance of all patched current classes before exit.
This is exactly the claimed front-end reduction to the one-corner best-seed
channel.
\end{proof}

\begin{corollary}[the clean replacement for the old $M1$ placeholder]
\label{cor:M1-clean}
In the rewritten one-cornered-odometer manuscript, the old broad front-end
placeholder should be replaced by the much sharper pair:
\begin{enumerate}[label=\textup{(\arabic*)},leftmargin=2.5em]
\item the endpoint-word shape theorem
      (Placeholder~\ref{ph:shape});
\item the initial entry theorem
      (Placeholder~\ref{ph:entry}).
\end{enumerate}
No additional post-entry non-trigger classification is needed in the main proof
spine.
\end{corollary}

\section{Bottom line}

The present note does not pretend to close the full front-end from raw
endpoint-search artifacts alone.
What it does prove is structurally important:
\begin{enumerate}[label=\textup{(\arabic*)},leftmargin=2.5em]
\item the normalized seed digits $[2,2,1]/[1,4,4]$ are just the normal form of
      an oriented bridge-type one-corner seed;
\item once the normalized seed enters the alternate-$2$ best-seed branch, the
      exact trigger family and all post-entry patched-class control are already
      theorem-level by the compact \textup{098} package;
\item therefore the true remaining front-end work is only a finite shape
      theorem plus a finite entry theorem.
\end{enumerate}

So the clean rewrite target is now precise:
odd-$m$, $d=5$ should be presented as a one-cornered odometer whose genuinely
new front-end burden is just the reduction to one normalized one-corner entry,
not a broad re-proving of the historical \textup{032}/\textup{033} artifact
chain.

\end{document}
